%% kapa: exam=2 round=37 year=2026 period=1 subject=appraisal_practice subject_ko=감정평가실무 seq=4 number=4 points=10
%% source: Q-Net 감정평가사 기출문제 (cst00309) artlSeq=5268128 file=raw/1008/5268128/26년 제37회 감정평가사 2차 (1교시_감정평가실무).pdf
%% conversion: kordoc 4.24.51 -> tools/gen2.py (auto); review=proofread-sonnet; flags=none
\begin{longquestion}{【문제 4】}
감정평가사 甲은 동일인 소유의 2개 필지에 대한 일반거래 목적의 감정평가를 의뢰받고, 자료를 수집하였다. 대상 토지는 「도시개발법」에서 규정하는 환지방식에 따른 사업시행지구 내에 있으며, 환지사항에 대한 조사내용은 아래의 자료와 같다. 다음 물음에 대하여 감정평가 방침을 약술하고, 감정평가액을 산정하기 위한 사정 면적을 확정하시오. (10점) \par
\medskip\noindent 물음 1) 환지처분 이전에 환지예정지로 지정된 후, 청산금이 납부되지 아니한 경우 (5점) \par
\medskip\noindent 물음 2) 환지예정지로 지정하기 이전인 경우 (5점) \par
< 자료 1 > 대상부동산 종전 자료 \par
1. 소재지: ○○시 ○○동 23번지, 24번지 \par
2. 지목: 23번지 ‘전’, 24번지 ‘답’ \par
3. 면적: 23번지 130m\textsuperscript{2}, 24번지 60m\textsuperscript{2} \par
4. 현황: 일단의 과수원 \par
5. 도시계획사항: 자연녹지지역 \par
< 자료 2 > 환지예정지 관련 수집 자료 \par
1. 종전토지: ○○시 ○○동 23번지, 24번지 \par
2. 환지예정지: ○○시 ○○동 4블록 5롯트 \par
3. 종전토지면적 합계: 190m\textsuperscript{2} \par
4. 권리면적: 160m\textsuperscript{2} \par
5. 환지(예정)면적: 180m\textsuperscript{2} \par
6. 징수면적: 20m\textsuperscript{2} \par
7. 도시계획(예정)사항: 제2종 일반주거지역 \par
\end{longquestion}
